% !TEX root = ../main-es.tex
% Las cajas de ayuda no son contenido: texcount las ignora (ver `make wordcount`).
%TC:macro \limite [ignore,ignore]
%TC:macro \nota [ignore]
%TC:macro \TODO [ignore]
%TC:macro \parencite [ignore]
%TC:macro \textcite [ignore]
\section{Posibles aplicaciones}
\limite{0}{Posibles aplicaciones: plantee posibles aplicaciones prácticas del tema
trabajado, señalando su utilidad en contextos reales.}

\begin{description}[leftmargin=1.4em]

\item[Diseño de motores de ejecución agéntica.] La aplicación directa se da en el motor que
motiva el proyecto, pues la frontera medida en $(\lambda, c)$ decide si conviene habilitar
la reparación en caliente, con qué presupuesto se reemplaza el máximo de iteraciones fijado
a mano y cuándo la gobernanza (un presupuesto $B$ de ediciones) supera a
la supervisión humana por aprobación. La misma lógica aplica a orquestadores de agentes LLM
genéricos, que hoy fijan por intuición cuándo replanificar.

\item[Presupuestos de deliberación con LLM.] Representar al LLM como un oráculo de
competencia $p$ y costo $c$ convierte una propiedad del modelo en un parámetro del
experimento: Las curvas resultantes dicen, para cualquier calidad y precio del modelo, si
la auto-edición de planes vale la pena. De esta forma un despliegue real calibra $p$ con
unas pocas corridas y lee su decisión en la frontera ya construida.

\item[Replanificación en robótica y control.] La pregunta de cuándo replanificar bajo
dinamismo y observabilidad parcial es central en control predictivo de POMDPs
\parencite{ma2018performance} y en la simplificación de lazo abierto con garantías
formales \parencite{kong2026open}, por lo que los resultados de este proyecto ofrecen una
versión medible y reproducible de ese compromiso que resulta útil como banco de pruebas
docente.

\item[Gobernanza de agentes que se modifican a sí mismos.] Un agente que reescribe su
propio programa con un alcance de edición acotado (una porción del plan que no puede
modificar) y un presupuesto de ediciones es un caso mínimamente
viable del problema general de alinear un agente automodificante: Qué parte de sí mismo
no debería poder tocar y cuántas veces conviene dejarlo intentarlo. Que el problema no es
hipotético lo documentaron los propios autores de The AI Scientist, cuyo agente editó su
código para extenderse el límite de tiempo que lo acotaba y en otra corrida se relanzó a sí
mismo \parencite{lu2024aiscientist}, de forma que la restricción escrita en el programa
resultó ser justamente lo que el programa podía reescribir. El agente gobernado de este
experimento mide qué tan costoso sale proteger esa frontera, pues el presupuesto $B$ (sumado
al alcance de edición acotado que comparten todas las arquitecturas que replanifican) es la versión medible de la pregunta por cuánta autoridad se le delega
a un agente sobre su propia especificación.

\end{description}
